% Chapter 8: Vector Fields
% Lee, Introduction to Smooth Manifolds (2nd ed., GTM 218).
% Generated from the section extraction; nodes carry only Lean that is
% verified sorry-free. See ../../UPSTREAM_LEAN_AUDIT.md.

\section{Vector Fields on Manifolds}

\begin{example}[Coordinate Vector Fields]
\label{ex:8-2}
\lean{smooth_chart_coordinate_vector_field_smooth}
\leanok


\dcref{ch8:8.2}
For a smooth chart $(U,(x^i))$ on $M$, the assignment
\[
p\longmapsto \left.\frac{\partial}{\partial x^i}\right|_p
\]
is a smooth vector field on $U$, called the $i$th coordinate vector field. Its
components in this chart are constant.
\end{example}

\begin{example}[The Euler Vector Field]
\label{ex:8-3}
\lean{euler_vector_field_points_radially_outward}
\leanok


\dcref{ch8:8.3}
The Euler vector field on $\mathbb{R}^n$ is
\[
V_x=\sum_{i=1}^n x^i\left.\frac{\partial}{\partial x^i}\right|_x.
\]
It is smooth, vanishes at the origin, and points radially outward elsewhere.
\end{example}

\begin{example}[The Angle Coordinate Vector Field on the Circle]
\label{ex:8-4}
\lean{exists_smooth_circle_angle_vector_field}
\leanok


\dcref{ch8:8.4}
If $\theta$ and $\widetilde\theta$ are angle coordinates on $\mathbb{S}^1$,
then locally $\widetilde\theta-\theta$ is constant, so
$d/d\theta=d/d\widetilde\theta$ on overlapping domains. These local fields
therefore define a global smooth vector field on $\mathbb{S}^1$, denoted
$d/d\theta$.
\end{example}

\begin{lemma}[Extension Lemma for Vector Fields]
\label{lem:8-6}
\lean{exists_supported_contMDiff_vectorField_extension_of_isClosed}
\leanok


\dcref{ch8:8.6}
Let $M$ be a smooth manifold, possibly with boundary, let $A\subseteq M$ be
closed, and let $X$ be a smooth vector field along $A$. For every open
$U\supseteq A$, there is a smooth vector field $\widetilde X$ on $M$ such that
$\widetilde X|_A=X$ and $\operatorname{supp}\widetilde X\subseteq U$.
\end{lemma}

\begin{proposition}
\label{prop:8-7}
\lean{exists_contMDiff_vectorField_eq_at_point}
\uses{lem:8-6}
\leanok



\dcref{ch8:8.7}
If $M$ is a smooth manifold, possibly with boundary, $p\in M$, and
$v\in T_pM$, then some smooth global vector field $X$ on $M$ satisfies
$X_p=v$.
\end{proposition}

\begin{proof}
\leanok
\sketch Extend $v$ to a constant-coefficient field in a coordinate
neighborhood of $p$, regard it as a smooth field along $\{p\}$, and apply the
extension lemma with $A=\{p\}$ and $U=M$.
\end{proof}

\begin{proposition}
\label{prop:8-8}
\lean{smoothVectorFieldModule}
\leanok


\dcref{ch8:8.8}
Let $M$ be a smooth manifold, possibly with boundary.
\begin{enumerate}
\item If $X,Y\in\mathfrak{X}(M)$ and $f,g\in C^\infty(M)$, then
$fX+gY$ is smooth.
\item Consequently, $\mathfrak{X}(M)$ is a module over $C^\infty(M)$.
\end{enumerate}
\end{proposition}


\section{Local and Global Frames}

\begin{definition}
\label{def:8-55-extra-1}
\lean{VectorField.SpansTangentBundleOn}
\leanok


An ordered tuple $(X_1,\ldots,X_k)$ of vector fields on $A\subseteq M$ is
\emph{linearly independent} if $(X_1|_p,\ldots,X_k|_p)$ is linearly
independent in $T_pM$ for every $p\in A$, and it
\emph{spans the tangent bundle} if these vectors span $T_pM$ for every $p$.
For an $n$-manifold, a \emph{local frame} on an open set $U$ is an ordered
$n$-tuple $(E_1,\ldots,E_n)$ whose values form a basis of every $T_pM$,
$p\in U$. It is \emph{global} when $U=M$ and \emph{smooth} when every $E_i$
is smooth. For an $n$-tuple, either pointwise independence or pointwise
spanning suffices to verify the frame condition.
\end{definition}

\begin{proposition}[Completion of Local Frames]
\label{prop:8-11}
\lean{exists_localFrame_extension_of_closed}
\uses{prob:8-5}
\leanok


\dcref{ch8:8.11}
Let $M$ be a smooth $n$-manifold, possibly with boundary.
\begin{enumerate}
\item If $(X_1,\ldots,X_k)$ is a linearly independent tuple of smooth vector
fields on an open $U\subseteq M$, where $1\leq k<n$, then each $p\in U$ has a
neighborhood $V$ carrying smooth fields $X_{k+1},\ldots,X_n$ such that
$(X_1,\ldots,X_n)$ is a smooth frame on $U\cap V$.
\item If $(v_1,\ldots,v_k)$ is linearly independent in $T_pM$, where
$1\leq k\leq n$, then a smooth local frame $(X_i)$ exists near $p$ with
$X_i|_p=v_i$ for $1\leq i\leq k$.
\item If $(X_1,\ldots,X_n)$ is a pointwise linearly independent tuple of
smooth vector fields along a closed $A\subseteq M$, then some neighborhood of
$A$ admits a smooth frame $(\widetilde X_1,\ldots,\widetilde X_n)$ satisfying
$\widetilde X_i|_A=X_i$ for every $i$.
\end{enumerate}
\end{proposition}

\begin{definition}
\label{def:8-55-extra-2}
\lean{IsOrthonormalFrameOn}
\leanok


A tuple $(E_1,\ldots,E_k)$ of vector fields on $A\subseteq\mathbb{R}^n$ is
\emph{orthonormal} if $(E_1|_p,\ldots,E_k|_p)$ is orthonormal for the Euclidean
inner product at every $p\in A$. An orthonormal frame is a frame consisting of
orthonormal vector fields.
\end{definition}

\begin{example}[A Polar Orthonormal Frame]
\label{ex:8-12}
\lean{example_8_12_E2_tangent_to_centered_circles}
\leanok


\dcref{ch8:8.12}
The standard coordinate frame is a global orthonormal frame on $\mathbb{R}^n$.
On $\mathbb{R}^2\setminus\{0\}$, with $r=\sqrt{x^2+y^2}$,
\[
E_1=\frac{x}{r}\frac{\partial}{\partial x}
   +\frac{y}{r}\frac{\partial}{\partial y},\qquad
E_2=-\frac{y}{r}\frac{\partial}{\partial x}
   +\frac{x}{r}\frac{\partial}{\partial y}. \tag{8.3}
\]
This is an orthonormal frame; $E_1$ is radial and $E_2$ is tangent to centered
circles.
\end{example}

\begin{lemma}[Gram--Schmidt Algorithm for Frames]
\label{lem:8-13}
\lean{exists_orthonormalFrameOn_with_same_initial_spans}
\leanok


\dcref{ch8:8.13}
If $(X_j)$ is a smooth local frame on an open $U\subseteq\mathbb{R}^n$, then
there is a smooth orthonormal frame $(E_j)$ on $U$ such that, for every
$p\in U$ and $1\leq j\leq n$,
\[
\operatorname{span}(E_1|_p,\ldots,E_j|_p)
=\operatorname{span}(X_1|_p,\ldots,X_j|_p).
\]
\end{lemma}

\begin{proof}
\leanok
\sketch Apply Gram--Schmidt pointwise:
\[
E_j=\frac{X_j-\sum_{i=1}^{j-1}(X_j\cdot E_i)E_i}
{\left|X_j-\sum_{i=1}^{j-1}(X_j\cdot E_i)E_i\right|}.
\]
The frame hypothesis makes every denominator nonzero. The formula therefore
defines smooth fields and preserves each initial span.
\end{proof}

\begin{definition}
\label{def:8-55-extra-3}
\lean{parallelizable}
\leanok


A smooth manifold, possibly with boundary, is \emph{parallelizable} if it
admits a smooth global frame.
\end{definition}

\section{Vector Fields as Derivations of $C^\infty(M)$}

\begin{notation}
\label{not:8-56-extra-1}
\lean{Notation_8_56_extra_1}
\leanok


For $X\in\mathfrak{X}(M)$ and a smooth $f\colon U\to\mathbb{R}$ on an open
$U\subseteq M$, define $Xf\colon U\to\mathbb{R}$ by
\[
(Xf)(p)=X_pf.
\]
This operation is local: for every open $V\subseteq U$,
\[
(Xf)|_V=X(f|_V). \tag{8.4}
\]
Here $fX$ denotes scalar multiplication of a vector field, whereas $Xf$
denotes the action of $X$ on a function.
\end{notation}

\section{Vector Fields and Smooth Maps}

\begin{definition}
\label{def:8-57-extra-1}
\lean{VectorField.f_related}
\leanok


Let $F\colon M\to N$ be smooth, and let $X$ and $Y$ be vector fields on $M$
and $N$, respectively. They are \emph{$F$-related} if
\[
dF_p(X_p)=Y_{F(p)}
\qquad\text{for every }p\in M.
\]
For a general smooth $F$, the vectors $dF_p(X_p)$ need not determine a vector
field on $N$: points outside $F(M)$ receive no value, and distinct points in a
fiber may yield incompatible values.
\end{definition}

\begin{proposition}
\label{prop:8-16}
\lean{f_related_iff_mfderiv_comp_eq}
\leanok


\dcref{ch8:8.16}
Let $F\colon M\to N$ be smooth between manifolds, possibly with boundary, and
let $X\in\mathfrak{X}(M)$ and $Y\in\mathfrak{X}(N)$. Then $X$ and $Y$ are
$F$-related if and only if every smooth real-valued function $f$ on an open
subset of $N$ satisfies
\[
X(f\circ F)=(Yf)\circ F. \tag{8.6}
\]
\end{proposition}

\begin{proof}
\leanok
\sketch At each $p\in M$ and for every smooth $f$ near $F(p)$,
\[
X(f\circ F)(p)=dF_p(X_p)f,
\qquad ((Yf)\circ F)(p)=Y_{F(p)}f.
\]
Equality for all such $f$ is equivalent to $dF_p(X_p)=Y_{F(p)}$.
\end{proof}

\begin{example}[The Rotation Field]
\label{ex:8-17}
\lean{example_8_17_d_dt_related_rotation_field}
\leanok


\dcref{ch8:8.17}
For $F\colon\mathbb{R}\to\mathbb{R}^2$ given by
$F(t)=(\cos t,\sin t)$, the field $d/dt$ is $F$-related to
\[
Y=x\frac{\partial}{\partial y}-y\frac{\partial}{\partial x}.
\]
\end{example}


\begin{proposition}
\label{prop:8-19}
\lean{existsUnique_smooth_f_related_vectorField_of_diffeomorph}
\leanok


\dcref{ch8:8.19}
If $F\colon M\to N$ is a diffeomorphism between smooth manifolds, possibly
with boundary, then every $X\in\mathfrak{X}(M)$ is $F$-related to a unique
smooth vector field on $N$.
\end{proposition}

\begin{proof}
\leanok
\sketch The unique possible field is
\[
Y_q=dF_{F^{-1}(q)}(X_{F^{-1}(q)}).
\]
It is smooth because, as a map $N\to TN$, it is the composite
\[
N\xrightarrow{F^{-1}}M\xrightarrow{X}TM\xrightarrow{dF}TN.
\]
\end{proof}

\begin{definition}
\label{def:8-57-extra-2}
\lean{VectorField.mpullback}
\uses{prop:8-19}

\mathlibok

For a diffeomorphism $F\colon M\to N$ and $X\in\mathfrak{X}(M)$, the unique
$F$-related field is the \emph{pushforward} $F_*X$, given by
\[
(F_*X)_q=dF_{F^{-1}(q)}(X_{F^{-1}(q)}). \tag{8.7}
\]
\end{definition}

\begin{corollary}
\label{cor:8-21}
\lean{pushforward_apply_comp_eq}
\uses{prop:8-16}
\leanok



\dcref{ch8:8.21}
If $F\colon M\to N$ is a diffeomorphism, $X\in\mathfrak{X}(M)$, and
$f\in C^\infty(N)$, then
\[
((F_*X)f)\circ F=X(f\circ F).
\]
\end{corollary}

\section{Vector Fields and Submanifolds}

\begin{definition}
\label{def:8-58-extra-1}
\lean{VectorField.IsTangentToSubmanifold}
\leanok


Let $S\subseteq M$ be an immersed or embedded submanifold, possibly with
boundary. A vector field $X$ on $M$ is \emph{tangent to $S$ at $p\in S$} if
$X_p\in T_pS\subseteq T_pM$, and it is \emph{tangent to $S$} if this holds at
every point of $S$.
\end{definition}

\section{Lie Brackets}

\begin{definition}
\label{def:8-59-extra-1}
\lean{VectorField.mlieBracket}

\mathlibok
For $X,Y\in\mathfrak{X}(M)$, their \emph{Lie bracket} is the operator
\[
[X,Y]f=XYf-YXf
\qquad (f\in C^\infty(M)).
\]
\end{definition}

\begin{proposition}[Coordinate Formula for the Lie Bracket]
\label{prop:8-26}
\lean{lieBracketWithin_coordinateComponent}
\leanok


\dcref{ch8:8.26}
Let $X$ and $Y$ be smooth vector fields on a smooth manifold $M$ with or
without boundary. In smooth local coordinates, if
$X=X^i\partial/\partial x^i$ and $Y=Y^j\partial/\partial x^j$, then
\[
[X,Y]=\left(X^i\frac{\partial Y^j}{\partial x^i}
-Y^i\frac{\partial X^j}{\partial x^i}\right)
\frac{\partial}{\partial x^j}, \tag{8.8}
\]
equivalently,
\[
[X,Y]=(XY^j-YX^j)\frac{\partial}{\partial x^j}. \tag{8.9}
\]
\end{proposition}

\begin{proof}
\leanok
\sketch Expanding $XYf-YXf$ in a chart gives first-derivative terms with the
displayed coefficients and second-derivative terms
$X^iY^j\partial_i\partial_jf-Y^jX^i\partial_j\partial_if$. The latter cancel
because mixed partial derivatives commute; relabeling dummy indices yields
(8.8).
\end{proof}

\begin{example}[A Lie Bracket on $\mathbb{R}^3$]
\label{ex:8-27}
\lean{example_8_27_lie_bracket_apply}
\leanok


\dcref{ch8:8.27}
For
\[
X=x\frac{\partial}{\partial x}+\frac{\partial}{\partial y}
  +x(y+1)\frac{\partial}{\partial z},
\qquad
Y=\frac{\partial}{\partial x}+y\frac{\partial}{\partial z},
\]
formula (8.9) gives
\[
[X,Y]=-\frac{\partial}{\partial x}-y\frac{\partial}{\partial z}.
\]
\end{example}

\begin{proposition}[Properties of the Lie Bracket]
\label{prop:8-28}
\lean{lie_bracket_smul_smul}
\leanok


\dcref{ch8:8.28}
For $X,Y,Z\in\mathfrak{X}(M)$, $a,b\in\mathbb{R}$, and
$f,g\in C^\infty(M)$, the Lie bracket satisfies
\begin{align*}
[aX+bY,Z]&=a[X,Z]+b[Y,Z],\\
[Z,aX+bY]&=a[Z,X]+b[Z,Y],\\
[X,Y]&=-[Y,X],\\
[X,[Y,Z]]+[Y,[Z,X]]+[Z,[X,Y]]&=0,\\
[fX,gY]&=fg[X,Y]+(fXg)Y-(gYf)X. \tag{8.11}
\end{align*}
\end{proposition}

\begin{proof}
\leanok
\sketch Bilinearity and antisymmetry follow directly from the definition.
Expanding the Jacobi sum on an arbitrary smooth function cancels all six
third-order compositions in pairs. Expanding $[fX,gY]$ and applying the
Leibniz rule gives (8.11).
\end{proof}


\begin{proposition}[Naturality of the Lie Bracket]
\label{prop:8-30}
\lean{f_related_mlieBracket}
\uses{prop:8-16}
\leanok



\dcref{ch8:8.30}
Let $F\colon M\to N$ be smooth between manifolds, possibly with boundary. If
$X_i\in\mathfrak{X}(M)$ is $F$-related to $Y_i\in\mathfrak{X}(N)$ for
$i=1,2$, then $[X_1,X_2]$ is $F$-related to $[Y_1,Y_2]$.
\end{proposition}

\begin{proof}
\leanok
\sketch Proposition~\ref{prop:8-16} gives, for every smooth $f$ on $N$,
\[
X_1X_2(f\circ F)=(Y_1Y_2f)\circ F,
\qquad
X_2X_1(f\circ F)=(Y_2Y_1f)\circ F.
\]
Subtracting and applying that proposition again proves the claim.
\end{proof}

\begin{corollary}[Pushforwards of Lie Brackets]
\label{cor:8-31}
\lean{pushforward_mlieBracket}
\uses{prop:8-30}
\leanok



\dcref{ch8:8.31}
If $F\colon M\to N$ is a diffeomorphism and
$X_1,X_2\in\mathfrak{X}(M)$, then
\[
F_*[X_1,X_2]=[F_*X_1,F_*X_2].
\]
\end{corollary}

\begin{proof}
\leanok
\sketch Apply Proposition~\ref{prop:8-30} with $Y_i=F_*X_i$ and use the
uniqueness of the pushforward.
\end{proof}

\section{The Lie Algebra of a Lie Group}

\begin{definition}
\label{def:8-60-extra-1}
\lean{VectorField.isLeftInvariant_iff_mfderiv}
\leanok


For a Lie group $G$, write $L_g(h)=gh$. A vector field $X$ on $G$ is
\emph{left-invariant} if it is $L_g$-related to itself for every $g\in G$;
equivalently,
\[
d(L_g)_{g'}(X_{g'})=X_{gg'}
\qquad(g,g'\in G). \tag{8.12}
\]
Since every $L_g$ is a diffeomorphism, this condition is also
$(L_g)_*X=X$. Smooth left-invariant fields form a linear subspace of
$\mathfrak{X}(G)$.
\end{definition}

\begin{proposition}
\label{prop:8-33}
\lean{isLeftInvariant_mlieBracket}
\uses{cor:8-31}
\leanok



\dcref{ch8:8.33}
The Lie bracket of two smooth left-invariant vector fields on a Lie group is
left-invariant.
\end{proposition}

\begin{proof}
\leanok
\sketch For every $g\in G$, Corollary~\ref{cor:8-31} gives
\[
(L_g)_*[X,Y]=[(L_g)_*X,(L_g)_*Y]=[X,Y].
\]
\end{proof}

\begin{definition}
\label{def:8-60-extra-2}
\lean{lie_jacobi}

\mathlibok
A \emph{Lie algebra} is a real vector space $\mathfrak g$ with a bilinear
bracket $[\ ,\ ]\colon\mathfrak g\times\mathfrak g\to\mathfrak g$ such that,
for all $X,Y,Z\in\mathfrak g$,
\[
[X,Y]=-[Y,X],
\qquad
[X,[Y,Z]]+[Y,[Z,X]]+[Z,[X,Y]]=0.
\]
All Lie algebras are taken over $\mathbb{R}$.
\end{definition}

\begin{definition}
\label{def:8-60-extra-3}
\lean{LieSubalgebra}

\mathlibok
A linear subspace $\mathfrak h\subseteq\mathfrak g$ is a \emph{Lie
subalgebra} if $[X,Y]\in\mathfrak h$ for all $X,Y\in\mathfrak h$; it then
inherits the restricted bracket.
\end{definition}

\begin{definition}
\label{def:8-60-extra-4}
\lean{LieHom}

\mathlibok
A linear map $A\colon\mathfrak g\to\mathfrak h$ is a \emph{Lie algebra
homomorphism} if
\[
A[X,Y]=[AX,AY].
\]
An invertible Lie algebra homomorphism is a \emph{Lie algebra isomorphism}.
\end{definition}

\begin{proposition}[Kernel and Image of a Lie Algebra Homomorphism]
\label{exer:8-34}
\lean{f.ker}

\mathlibok
\dcref{ch8:8.34}
The kernel and image of a Lie algebra homomorphism are Lie subalgebras.
\end{proposition}

\begin{proposition}[Basis Criterion for Lie Algebra Homomorphisms]
\label{exer:8-35}
\lean{linearMap_is_lieHom_iff_preserves_bracket_on_basis}
\leanok


\dcref{ch8:8.35}
Let $A\colon\mathfrak g\to\mathfrak h$ be linear between finite-dimensional
Lie algebras. Then $A$ is a Lie algebra homomorphism if and only if, for
some basis $(E_1,\ldots,E_n)$ of $\mathfrak g$,
\[
A[E_i,E_j]=[AE_i,AE_j]
\qquad(1\leq i,j\leq n).
\]
\end{proposition}

\begin{example}[Standard Lie Algebras]
\label{ex:8-36}
\lean{abelian_lie_algebra}
\uses{prop:8-28}
\leanok



\dcref{ch8:8.36}
The following are Lie algebras or Lie subalgebras.
\begin{enumerate}
\item $\mathfrak{X}(M)$ with the vector-field bracket.
\item The smooth left-invariant vector fields on a Lie group $G$.
\item The $n^2$-dimensional space
$\mathfrak{gl}(n,\mathbb{R})=\mathrm{M}(n,\mathbb{R})$ with
$[A,B]=AB-BA$.
\item The real $2n^2$-dimensional space
$\mathfrak{gl}(n,\mathbb{C})=\mathrm{M}(n,\mathbb{C})$ with the same bracket.
\item For a vector space $V$, the space $\mathfrak{gl}(V)$ of endomorphisms,
with $[A,B]=A\circ B-B\circ A$; under the standard matrix identification,
$\mathfrak{gl}(\mathbb R^n)=\mathfrak{gl}(n,\mathbb R)$.
\item Any vector space with the zero bracket; such a Lie algebra is
\emph{abelian}.
\end{enumerate}
\end{example}

\begin{definition}
\label{def:8-60-extra-5}
\lean{contMDiff_mulInvariantVectorField}

\mathlibok
The \emph{Lie algebra of a Lie group $G$}, denoted $\operatorname{Lie}(G)$,
is the Lie algebra of its smooth left-invariant vector fields.
\end{definition}

\begin{theorem}
\label{thm:8-37}
\lean{lie_algebra_finrank_eq_manifold_finrank}
\leanok


\dcref{ch8:8.37}
For a Lie group $G$, evaluation at the identity,
\[
\varepsilon\colon\operatorname{Lie}(G)\longrightarrow T_eG,
\qquad \varepsilon(X)=X_e,
\]
is a vector-space isomorphism. Hence
$\dim\operatorname{Lie}(G)=\dim G$.
\end{theorem}

\begin{proof}
\leanok
\sketch The map is linear. Left invariance implies
$X_g=d(L_g)_e(X_e)$, so evaluation is injective. For $v\in T_eG$, define
\[
v^L|_g=d(L_g)_e(v). \tag{8.13}
\]
If $\gamma(0)=e$ and $\gamma'(0)=v$, then for $f\in C^\infty(G)$,
\[
(v^Lf)(g)=\left.\frac{d}{dt}\right|_{t=0}f(g\gamma(t)).
\]
Smooth dependence on $(t,g)$ shows that $v^L$ is smooth. The identity
$L_h\circ L_g=L_{hg}$ makes $v^L$ left-invariant, and $v^L|_e=v$; thus
evaluation is surjective.
\end{proof}

\begin{notation}
\label{not:8-60-extra-6}
\lean{mulInvariantVectorField}

\mathlibok
For $v\in T_eG$, write $v^L$ for the smooth left-invariant vector field
defined by (8.13).
\end{notation}

\begin{corollary}
\label{cor:8-38}
\lean{left_invariant_rough_vector_field_smooth}
\leanok


\dcref{ch8:8.38}
Every left-invariant rough vector field on a Lie group is smooth.
\end{corollary}

\begin{proof}
\leanok
\sketch If $v=X_e$, left invariance gives $X=v^L$, which is smooth by the
theorem.
\end{proof}

\begin{definition}
\label{def:8-60-extra-7}
\lean{IsLeftInvariantFrameOn}
\leanok


A local or global frame on a Lie group is \emph{left-invariant} if each of its
fields is left-invariant.
\end{definition}

\begin{corollary}
\label{cor:8-39}
\lean{lie_group_is_parallelizable}
\leanok


\dcref{ch8:8.39}
Every Lie group has a smooth left-invariant global frame and is therefore
parallelizable.
\end{corollary}

\begin{proof}
\leanok
\sketch Under the evaluation isomorphism, any basis of $T_eG$ corresponds to
a basis of $\operatorname{Lie}(G)$ whose pointwise values form a global frame.
\end{proof}

\begin{proposition}[Lie Algebra of the General Linear Group]
\label{prop:8-41}
\lean{general_linear_group_lie_equiv_matrix}
\leanok


\dcref{ch8:8.41}
The natural maps
\[
\operatorname{Lie}(\mathrm{GL}(n,\mathbb{R}))
\longrightarrow T_{I_n}\mathrm{GL}(n,\mathbb{R})
\longrightarrow\mathfrak{gl}(n,\mathbb{R}) \tag{8.14}
\]
compose to a Lie algebra isomorphism.
\end{proposition}

\begin{proof}
\leanok
\sketch Under the coordinates $X_j^i$ and the identification
\[
\left.A_j^i\frac{\partial}{\partial X_j^i}\right|_{I_n}
\longleftrightarrow (A_j^i),
\]
where $i,j$ serve simultaneously as coordinate and matrix row/column indices;
the summation convention applies, and an index in a denominator is interpreted
with the opposite variance.
The left-invariant field associated with $A\in\mathfrak{gl}(n,\mathbb{R})$ is
\[
A^L|_X=\left.X_j^iA_k^j\frac{\partial}{\partial X_k^i}\right|_X. \tag{8.15}
\]
The coordinate bracket formula yields
\[
[A^L,B^L]|_{I_n}
=\left(A_k^iB_r^k-B_k^iA_r^k\right)
 \left.\frac{\partial}{\partial X_r^i}\right|_{I_n},
\]
which corresponds to the matrix commutator $AB-BA$. Since a left-invariant
field is determined by its value at $I_n$, $[A^L,B^L]=[A,B]^L$.
\end{proof}

\section{The Lie Algebra of a Lie Subgroup}

\begin{proposition}[The Lie Algebra of $\mathrm{GL}(n,\mathbb{C})$]
\label{prop:8-48}
\lean{complex_general_linear_group_identity_tangent_equiv_matrix_apply}
\uses{ex:7-18, prop:8-41}
\leanok



\dcref{ch8:8.48}
Let $\varepsilon$ be evaluation at the identity and let $\varphi$ identify the
tangent space of the open matrix group with the matrix algebra. The maps
\[
\operatorname{Lie}(\mathrm{GL}(n,\mathbb{C}))
\xrightarrow{\ \varepsilon\ }
T_{I_n}\mathrm{GL}(n,\mathbb{C})
\xrightarrow{\ \varphi\ }
\mathfrak{gl}(n,\mathbb{C}) \tag{8.18}
\]
compose to a Lie algebra isomorphism.
\end{proposition}

\begin{proof}
\leanok
\sketch The real representation
$\beta\colon\mathrm{GL}(n,\mathbb{C})\to\mathrm{GL}(2n,\mathbb{R})$ induces
$\beta_* $ on Lie algebras. Because $\beta$ is linear, its derivative at
$I_n$ has the same formula. Under the canonical tangent-space identifications
this gives an injective map
\[
\alpha\colon\mathfrak{gl}(n,\mathbb{C})
\longrightarrow\mathfrak{gl}(2n,\mathbb{R}).
\]
The identity $\beta(AB)=\beta(A)\beta(B)$ implies
\[
\alpha([A,B])=[\alpha(A),\alpha(B)].
\]
The commutative diagram with Proposition~\ref{prop:8-41}, restricted to the
images of the vertical maps, therefore shows that the maps in (8.18) preserve
brackets and form the asserted isomorphism.
\end{proof}

\begin{definition}
\label{def:8-62-extra-1}
\lean{LieModule.IsFaithful}

\mathlibok
A finite-dimensional \emph{representation} of a finite-dimensional Lie
algebra $\mathfrak g$ is a Lie algebra homomorphism
$\varphi\colon\mathfrak g\to\mathfrak{gl}(V)$ for a finite-dimensional vector
space $V$. It is \emph{faithful} if $\varphi$ is injective; then
$\mathfrak g\cong\varphi(\mathfrak g)\subseteq\mathfrak{gl}(V)$.
\end{definition}

\begin{exercise}
\label{exer:8-18}
\lean{example_8_17_d_dt_related_rotation_field}
\leanok
\uses{ex:8-17,prop:8-16}
\dcref{ch8:8.18}
Prove the claim in the preceding example in two ways: directly from the definition, and by using Proposition 8.16.
\end{exercise}

\begin{exercise}
\label{exer:8-29}
\lean{lie_bracket_smul_smul}
\leanok
\uses{prop:8-28}
\dcref{ch8:8.29}
Prove part (d) of the preceding proposition.
\end{exercise}

\begin{exercise}
\label{exer:8-9}
\lean{smoothVectorFieldModule}
\leanok
\uses{prop:8-8}
\dcref{ch8:8.9}
\begin{itemize}
\item Exercise 8.9. Prove Proposition 8.8.
\end{itemize}
\end{exercise}
